% Compsoc guidelines: 100~200 word limit. Must clearly state the nature and significance of the paper; must not include mathematical expressions or bibliographic references.
\IEEEtitleabstractindextext{
  \begin{abstract}
    Accelerating 3D CNNs poses a significant design challenge, due in large part to the increase in arithmetic complexity of the convolution operation. Winograd-based designs can achieve high resource efficiency in networks using small kernels, yet are limited by increasingly complex transform operations in higher dimensions. FFT convolution has seen recent success in applying overlap-add to CNNs, increasing its efficiency for large inputs. Likewise, the Complex-domain multiplications required by the algorithm can be optimized using polar coordinate representation. We propose a 3D CNN accelerator using FFT convolution augmented by overlap-add and CORDIC algorithms. Our analysis reveals the lower assymptotic bound of overlap-add FFT compared to Winograd convolution. We implement and test this design on FPGA technology, achieving over double DSP-efficiency compared to the state-of-the-art.
  \end{abstract}
  \begin{IEEEkeywords}
    % COMPSOC: http://www.computer.org/mc/keywords/keywords.htm
    % IEEE: https://ieee-org.widen.net/s/jwk9pcxxvd/ieee-taxonomy
    Convolutional neural networks,
    % Neural network hardware,
    % Open source hardware,
    Fast Fourier transforms,
    Computational efficiency,
    Algorithmic efficiency,
    Field programmable gate arrays,
    % AI accelerators,
    Accelerator architectures,
    Hardware acceleration.
  \end{IEEEkeywords}
}